\section{引言}
\label{sec:intro}

部分子物理对于理解强子高能碰撞的动力学以及研究强子的内部结构都十分重要~\cite{Ellis:1991qj,Thomas:2001kw}。最熟悉的例子是夸克与胶子的部分子分布函数（PDF）：它们一方面为对撞机上的高能产生过程提供束流信息~\cite{Gao:2017yyd}，另一方面描述了对撞强子的束缚态物理。

尽管十分重要，从量子色动力学（QCD）的第一性原理计算部分子物理一直是个挑战。最近，人们提出了一种有效场论（EFT）方法——大动量有效理论（LaMET）——从以大动量运动的强子的物理性质中提取部分子物理~\cite{Ji:2013dva,Ji:2014gla}；后者可以通过对欧氏 QCD 的系统近似（如格点场论）来计算。自提出以来，LaMET 已被广泛用于计算夸克同位旋矢量分布函数~\cite{Lin:2014zya,Alexandrou:2015rja,Chen:2016utp,Alexandrou:2016jqi,Alexandrou:2018pbm,Chen:2018xof,Lin:2018qky,Liu:2018uuj,Alexandrou:2018eet,Liu:2018hxv,Chen:2018fwa,Izubuchi:2019lyk,Shugert:2020tgq,Chai:2020nxw,Lin:2020ssv,Fan:2020nzz}、分布振幅（DA）~\cite{Zhang:2017bzy,Chen:2017gck,Zhang:2020gaj}、广义部分子分布~\cite{Chen:2019lcm,Alexandrou:2019dax}，以及近来的横动量依赖分布~\cite{Shanahan:2019zcq,Shanahan:2020zxr,Zhang:2020dbb}，甚至高扭度分布~\cite{Bhattacharya:2020cen}。关于 LaMET 的一些近期综述可参见文献~\cite{Ji:2020ect,Cichy:2018mum}。

LaMET 的核心思想是：无穷大动量系（IMF）中的部分子可以用具有较大但有限动量的强子的物理性质来近似。由于紫外（UV）发散的存在，这一近似并非完全直接明了，需要使用匹配（matching）与跑动（running）这类标准 EFT 技术。详细的研究表明，标准的 DGLAP 演化~\cite{Dokshitzer:1977sg,Gribov:1972ri,Altarelli:1977zs} 正源于强子物理性质随动量的演化~\cite{Ji:2020ect}。

在 LaMET 的应用中，人们从空间关联函数的格点计算出发。由于格点破坏连续时空对称性，裸关联函数中会出现幂次发散。在与连续方案（如维数正规化和（修正的）最小减除 ${\overline {\rm MS}}$）中的关联匹配时必须将其减除。过去在实际应用中提出的主要方法包括正规化无关动量减除法（RI/MOM）~\cite{Constantinou:2017sej,Stewart:2017tvs,Alexandrou:2017huk,Chen:2017mzz} 与比值方法~\cite{Radyushkin:2017cyf,Orginos:2017kos,Braun:2018brg,Li:2020xml}。后者依赖于欧氏算符乘积展开（OPE）的有效性，只能用于短距离关联，因而不能直接用于 LaMET 应用。相比之下，RI/MOM 方法乍看之下似乎适用于大 $z$（$z$ 为规范链长度）距离。然而细致考察表明，该方法会引入潜在的非微扰效应，例如通过方案匹配中形如 $\ln(z^2\mu^2)$（$\mu$ 为重正化标度）的红外（IR）对数。既然在 QCD 这类渐近自由理论中 UV 发散应当是微扰的，理应可以找到一种不引入非微扰效应的重正化程序。就目前格点 QCD 计算的精度而言，RI/MOM 或许并未引入大的非微扰效应；然而高精度的系统性有效场论计算无法回避这一问题。

为实现这一点，我们在本文中针对 LaMET 应用中的格点关联提出一种混合重正化方案。在 OPE 成立的短距离处，推荐采用标准的 RI/MOM 或比值方法；在大距离处，我们建议使用辅助场表述~\cite{Samuel:1978iy,Gervais:1979fv,Arefeva:1980zd,Dorn:1986dt}，该表述已被多位作者在 LaMET 应用中所倡导~\cite{Ji:2017oey,Green:2017xeu,Green:2020xco}。在此表述中，Wilson 线被辅助场的两点函数所取代：格点关联函数中的线性发散于是与辅助场的质量重正化相联系，而对数发散则出现在 Wilson 线端点处``重''-轻``流''的重正化中。两种发散都可以按与 $\overline{\rm MS}$ 方案一致的方式分别重正化~\cite{Green:2017xeu,
Green:2020xco}。尽管 Wilson 线的质量减除早有建议~\cite{Chen:2016fxx,Ishikawa:2017faj}，但据我们所知，由于文献中尚未确立一种计算非微扰质量的可靠途径，它一直未被广泛付诸实用。这里我们提出几种可行的做法，留待未来通过系统的格点模拟加以研究。

此外，我们还讨论了利用 LaMET 提取部分子物理时十分重要的其他几个问题，例如如何在 $z\sim 0$ 附近恰当地与连续方案匹配，以及如何利用相关关联函数在大光前（LF）距离处的渐近行为，消除截断傅里叶变换在动量分布中引起的非物理振荡。
